% ============================================================================
% Fichier  : partie1/P01_C01_pourquoi_superviser.tex
% Titre    : Pourquoi Superviser
% Auteur   : MINKA MI NGUIDJOÏ Thierry Emmanuel
% ============================================================================
% Sources : ISO/IEC 27035 ; NIST SP 800-61r2 ; CIS Controls v8. Opération
% BÉNOUÉ (LiZbiZ) : cas fictif.
% ============================================================================

\chapter{Pourquoi Superviser}
\label{chap:pourquoi}
\chaptermark{Pourquoi superviser}

\epigraph{%
    « On ne défend pas ce que l'on ne voit pas. »%
}{Principe du cours}

\niveauChapitre{Fondamental}{Le prologue~; notions de réseaux et de systèmes.}

\begin{boxobjectifs}
\begin{itemize}
    \item Définir la supervision de sécurité et la distinguer de la
          supervision d'exploitation.
    \item Situer la détection dans la chaîne prévention, détection, réaction.
    \item Comprendre le délai de présence de l'attaquant et pourquoi il
          compte.
    \item Relier la supervision aux obligations de l'entreprise.
    \item Poser les indicateurs fondamentaux de la détection.
\end{itemize}
\end{boxobjectifs}

\bigskip

% ============================================================================
\section*{Question directrice : si tout était dans les journaux, pourquoi personne n'a-t-il rien vu~?}
% ============================================================================

Onze jours de présence, des traces à chaque étape, et pourtant la découverte
par un écran de rançon. Ce chapitre explique pourquoi la prévention ne suffit
pas, ce qu'est superviser, et comment se mesure la capacité à voir venir.

\medskip
\textbf{Plan du chapitre.}
\begin{enumerate}
    \item Prévenir ne suffit pas (section~\ref{sec:c01-prevenir}).
    \item Qu'est-ce que superviser (section~\ref{sec:c01-definir}).
    \item Le délai de présence (section~\ref{sec:c01-delai}).
    \item Superviser, une obligation (section~\ref{sec:c01-obligation}).
    \item Les indicateurs fondamentaux (section~\ref{sec:c01-indicateurs}).
\end{enumerate}

\separateur

% ============================================================================
\section{Prévenir ne Suffit Pas}
\label{sec:c01-prevenir}
% ============================================================================

La sécurité se joue sur trois temps~: \termecle{prévention} (empêcher),
\termecle{détection} (voir ce qui est passé malgré la prévention) et
\termecle{réaction} (limiter et corriger). Aucune prévention n'est parfaite~:
un correctif manquant, un mot de passe faible, un courriel piégé suffisent.
La détection est le filet qui rattrape ce que la prévention a laissé passer.

% ============================================================================
\section{Qu'est-ce que Superviser}
\label{sec:c01-definir}
% ============================================================================

\begin{boxdef}
\termecle{Supervision de sécurité}~: collecte continue des traces d'un
système, leur corrélation et leur analyse pour détecter, le plus tôt
possible, les activités malveillantes ou anormales. Elle se distingue de la
\termecle{supervision d'exploitation}, qui surveille la disponibilité et la
performance, non les attaques.
\end{boxdef}

% ============================================================================
\section{Le Délai de Présence}
\label{sec:c01-delai}
% ============================================================================

\begin{boxdef}
\termecle{Délai de présence} (temps de séjour)~: durée entre la première
compromission et sa détection. Plus il est long, plus l'attaquant s'étend,
exfiltre et s'installe.
\end{boxdef}

\begin{boxscenario}{Opération BÉNOUÉ. Onze jours}
\textbf{Constatation.} Première trace le 23~février (connexion d'un compte de
service à une heure inhabituelle)~; chiffrement dans la nuit du 5 au 6 mars. Délai de présence~:
onze jours.

\medskip
\textbf{Analyse} (confiance élevée). Réduire ce délai de onze jours à
quelques heures aurait permis de contenir l'attaque avant l'exfiltration et
le chiffrement. C'est l'objectif premier de la supervision.
\end{boxscenario}

% ============================================================================
\section{Superviser, une Obligation}
\label{sec:c01-obligation}
% ============================================================================

\begin{boxjuris}
Au Cameroun, la loi n°~2010/012 du 21~décembre~2010 relative à la
cybersécurité et à la cybercriminalité impose la conservation des traces de
connexion et la coopération avec l'\sigle{ANTIC} en cas d'incident
(Chapitre~\ref{chap:droit}). Ne pas disposer de journaux exploitables n'est
pas seulement une faiblesse technique~: c'est un manquement opposable.
\end{boxjuris}

% ============================================================================
\section{Les Indicateurs Fondamentaux}
\label{sec:c01-indicateurs}
% ============================================================================

\begin{tableaucours}{L{4.4cm}L{7.8cm}}{%
    \tableauHeader{Indicateur} &
    \tableauHeader{Ce qu'il mesure}
}
Délai moyen de détection & Temps entre compromission et alerte \\
Délai moyen de réaction & Temps entre alerte et confinement \\
Couverture de journalisation & Part des systèmes dont les traces sont collectées \\
Taux de faux positifs & Alertes sans incident réel \\
Taux de détection & Part des attaques (ou techniques) effectivement vues \\
\end{tableaucours}
\vspace{-0.8em}
\begin{center}{\small\itshape Tableau~1.1. Indicateurs fondamentaux de la supervision}\end{center}

\begin{boxmetier}{Analyste SOC}
\textbf{Mesurer pour progresser.} L'analyste ne se contente pas de traiter
les alertes~: il suit le délai de détection d'un mois sur l'autre et
l'améliore. Un SOC qui ne mesure pas son délai de détection ne sait pas s'il
s'améliore.
\end{boxmetier}

\separateur

\begin{boxresume}
\textbf{Ce qu'il faut retenir de ce chapitre~:}
\begin{itemize}
    \item \textbf{Prévention, détection, réaction}~: la détection rattrape ce
          que la prévention laisse passer.
    \item \textbf{Superviser} la sécurité n'est pas superviser l'exploitation.
    \item \textbf{Le délai de présence} est l'ennemi~; le réduire est
          l'objectif.
    \item \textbf{Superviser est une obligation} (loi 2010/012).
    \item \textbf{On mesure} délai de détection, de réaction, couverture.
\end{itemize}
\end{boxresume}

\begin{boxinfo}
\textbf{Points de contrôle de l'élève.} Je sais~: définir la supervision de
sécurité~; situer la détection dans la chaîne~; expliquer le délai de
présence~; rappeler l'obligation légale~; citer les indicateurs.
\end{boxinfo}

\bigskip

\begin{boxexo}[title={\ding{45}\quad Exercice~1.1. Trois temps \niveaubase}]
Classez~: installer un correctif~; lire les journaux chaque matin~; isoler un
poste infecté. À quel temps (prévention, détection, réaction) chacun
appartient-il~?
\end{boxexo}

\begin{boxexo}[title={\ding{45}\quad Exercice~1.2. Deux supervisions \niveaubase}]
Distinguez supervision de sécurité et supervision d'exploitation par un
exemple de chacune.
\end{boxexo}

\begin{boxexo}[title={\ding{45}\quad Exercice~1.3. Le délai \niveaubase}]
Calculez le délai de présence de l'opération BÉNOUÉ. Pourquoi le réduire~?
\end{boxexo}

\begin{boxexo}[title={\ding{45}\quad Exercice~1.4. Le filet \niveaumoyen}]
« Puisque nous avons un pare-feu et un antivirus, la détection est inutile. »
Répondez.
\end{boxexo}

\begin{boxexo}[title={\ding{45}\quad Exercice~1.5. Obligation \niveaumoyen}]
En quoi l'absence de journaux exploitables est-elle un manquement opposable
chez LiZbiZ~?
\end{boxexo}

\begin{boxexo}[title={\ding{45}\quad Exercice~1.6. Indicateurs \niveaumoyen}]
Pour LiZbiZ, proposez une valeur cible de délai de détection et justifiez-la.
\end{boxexo}

\begin{boxexo}[title={\ding{45}\quad Exercice~1.7. Couverture \niveaumoyen}]
Pourquoi une couverture de journalisation partielle crée-t-elle des angles
morts~?
\end{boxexo}

\begin{boxexo}[title={\ding{45}\quad Exercice~1.8. Faux positifs \niveaumoyen}]
Quel risque un taux de faux positifs trop élevé fait-il courir à la
détection~?
\end{boxexo}

\begin{boxexo}[title={\ding{45}\quad Exercice~1.9. Tableau de bord \niveauexpert}]
Construisez le tableau de bord de supervision que LiZbiZ devrait suivre après
l'opération GAROUA.
\end{boxexo}

\begin{boxexo}[title={\ding{45}\quad Exercice~1.10. Note à la direction \niveauexpert}]
Rédigez une note qui explique à M.~ETEME pourquoi investir dans la détection,
et non seulement dans la prévention.
\end{boxexo}

\bigskip

\begin{boxinfo}
\textbf{Pour aller plus loin.}
\begin{itemize}
    \item ISO/IEC 27035~\cite{iso27035}~; NIST SP~800-61r2~\cite{nist80061r2}.
\end{itemize}
\end{boxinfo}

\begin{boxinfo}
\textbf{Transition.} Superviser suppose une organisation~: le
Chapitre~\ref{chap:soc} présente l'architecture d'un centre d'opérations de
sécurité.
\end{boxinfo}
